\section{Operational Engineering and Deployment Specification}
\label{sec:operational_engineering}

Deploying agentic reasoning architectures in production requires strict rate-limiting, environment isolation, and reproducible execution pipelines. This section details operational configurations, command-line interfaces, and deployment guidelines for MedRAGAgents.

\subsection{CLI Command Interface (\texttt{run.py})}

The main entry point \texttt{run.py} controls pipeline execution via command-line arguments. Table~\ref{tab:cli_params} details the available CLI parameters.

\begin{table}[h!]
\centering
\caption{Command-Line Interface Parameter Specifications for \texttt{run.py}.}
\label{tab:cli_params}
\small
\begin{tabularx}{\textwidth}{l c c X}
\toprule
\textbf{CLI Parameter} & \textbf{Data Type} & \textbf{Default Value} & \textbf{Operational Function} \\
\midrule
\texttt{--variant} & String & \texttt{V3} & Selects target system variant (\texttt{V0}, \texttt{V1}, \texttt{V2}, \texttt{V3}, or \texttt{ALL}). \\
\texttt{--n} & Integer & \texttt{-1} & Sets evaluation sample limit ($10$ for quick testing, $-1$ for full dataset). \\
\texttt{--start} & Integer & \texttt{0} & Specifies starting vignette index for parallel batch processing. \\
\texttt{--llm} & String & \texttt{cfg.DEFAULT\_LLM} & Sets foundation model backbone (\texttt{google/gemini-3.6-flash}). \\
\texttt{--delay} & Float & \texttt{2.0} & Enforces inter-vignette delay (seconds) to prevent API rate-limit errors. \\
\texttt{--dataset\_dir} & Path & \texttt{./datasets/MedQA/} & Specifies directory path to input \texttt{test.jsonl} file. \\
\texttt{--output\_dir} & Path & \texttt{./outputs} & Specifies output directory for generated prediction files. \\
\texttt{--evaluate} & Flag & \texttt{False} & Triggers automated evaluation script (\texttt{evaluate.py}) upon run completion. \\
\bottomrule
\end{tabularx}
\end{table}

\subsection{Production Execution Guidelines}

To execute a complete evaluation across all four variants on the MedQA-USMLE test set with automated statistical reporting:

\begin{lstlisting}[language=bash, caption={Production Execution Shell Command.}, label={lst:cli_run}]
# Step 1: Initialize Virtual Environment and Dependencies
python3 -m venv .venv
source .venv/bin/activate
pip install -r requirements.txt

# Step 2: Download MedCPT Textbooks Corpus Vector Index
python3 download_corpus.py

# Step 3: Execute All Variants with Rate-Limit Protection and Auto-Evaluation
PYTHONUNBUFFERED=1 python3 run.py \
  --variant ALL \
  --n 10 \
  --delay 2.0 \
  --evaluate 2>&1 | tee outputs/run_execution.log
\end{lstlisting}

\subsection{Rate-Limiting Buffer and Security Controls}

API requests are throttled using an adjustable inter-call delay parameter (\texttt{--delay}). On free-tier API quotas (e.g., 5 RPM limit), setting \texttt{--delay 12.0} maintains compliance with remote quota restrictions.

Environment variables are managed securely via \texttt{.env} files, isolating secret API credentials from the repository codebase:

\begin{lstlisting}[language=bash, caption={Environment Configuration File (.env).}, label={lst:env_spec}]
GEMINI_API_KEY=your_secure_api_key_here
DEFAULT_LLM=google/gemini-3.6-flash
DATASET_DIR=./datasets/MedQA/
OUTPUT_DIR=./outputs
LONG_TERM_MEM=./memory/long_term_cache.json
\end{lstlisting}
